\section{Proof of the main theorem}\label{sec:proof}

\begin{lemma}[Dodziuk]\label{lem:dodziuk}
Let $g,\bar g$ be smooth metrics on a closed $n$-manifold with
$(1-\delta)\bar g\le g\le(1+\delta)\bar g$, $0<\delta<1/2$. Then the
ordered positive eigenvalues of the Hodge Laplacians in each degree $p$
satisfy
\[
 \Theta(\delta)^{-1}\lambda^+_{p,j}(\bar g)\le\lambda^+_{p,j}(g)
 \le\Theta(\delta)\,\lambda^+_{p,j}(\bar g),\qquad
 \Theta(\delta)=1+O_n(\delta).
\]
\end{lemma}
\begin{proof}
This follows from \cite[Proposition~3.3]{dodziuk}. The positive spectrum in degree
$p$ is the union of the squared singular values of $d$ on
$\operatorname{Dom}d/\ker d$ in degrees $p$ and $p-1$, with the quotient
$L^2$ norms; these norms and the norms of the images are changed by
factors $1+O_n(\delta)$, and neither the operator $d$ nor its domain
depends on the metric.
\end{proof}

\begin{proof}[Proof of Theorem~\ref{thm:main}]
Let $t_l\to0$ be a sequence, and write $t_l=\epsilon_l^2u_l^2$ as in
\eqref{eq:scale}. Passing to a subsequence, $u_l\to u_0$. Fix $R$ and
$A$. By Lemma~\ref{lem:condition-N}, $X_{t_l}$ is identified with
$\overline X_{\epsilon_l,u_0}$ in such a way that $g_{t_l}$ and
$\bar g^{R,A}_{\epsilon_l,u_0}$ differ, relative to the latter, by at
most $\delta_l$ with $\limsup_l\delta_l\le\delta(R,A):=C(R+A^{-2})$.
Both manifolds have $s$ zero eigenvalues in degree two and none in degree
one. By Lemma~\ref{lem:dodziuk}, Proposition~\ref{prop:count} and
Theorem~\ref{thm:model-law}, if $\delta(R,A)<1/4$ then for large $l$
exactly $m$ positive degree-two eigenvalues of $g_{t_l}$ are small, the
degree-one spectrum and the other positive degree-two eigenvalues are
bounded below by a positive constant, and
\begin{align*}
 \Theta(2\delta(R,A))^{-1}c\,\mu_j(R)
 &\le\liminf_l\log(1/\epsilon_l)\lambda^+_{2,j}(t_l)\\
 &\le\limsup_l\log(1/\epsilon_l)\lambda^+_{2,j}(t_l)
 \le\Theta(2\delta(R,A))\,c\,\mu_j(R).
\end{align*}
Let $R\to0$ and $A\to\infty$ and use \eqref{eq:mass-limit}: along the
subsequence, $\log(1/\epsilon_l)\lambda^+_{2,j}(t_l)\to c\,\mu_j=\pi\mu_j$.
Since $\log(1/\epsilon)=\tfrac12\log(1/|t|)$ this is \eqref{eq:main}
along the subsequence. As every sequence has such a subsequence,
\eqref{eq:main} holds, and so do the count and the uniform lower bound
for the remaining spectrum.

Let $\psi$ be an eigenform for a small eigenvalue $\lambda$, with $|t|$
so small that $\lambda$ is below the bound $\gamma$ for the degree-one
spectrum. Then $d^*\psi$ is a one-form with $\Delta d^*\psi=\lambda
d^*\psi$, hence zero, and $\psi=\lambda^{-1}d^*d\psi$ is coexact. Since
$g_t$ is K\"ahler, $\Lambda$ commutes with $\Delta$, so $\Lambda\psi$ is
a function with $\Delta\Lambda\psi=\lambda\Lambda\psi$. The positive
spectrum on functions is the spectrum on exact one-forms and is at least
$\gamma$; hence $\Lambda\psi=0$. The type is the subject of the next
proposition.
\end{proof}

\begin{proposition}\label{prop:type}
For $|t|$ sufficiently small, every eigenform of
$(X_t,g_t)$ for one of the eigenvalues
$\lambda^+_{2,1}(t),\ldots,\lambda^+_{2,m}(t)$ is of type $(1,1)$.
\end{proposition}
\begin{proof}
Let $S_t$ be the span of these eigenforms and $P_t$ the pointwise
orthogonal projection of two-forms on $X_t$ onto the forms of type
$(2,0)+(0,2)$. Since $g_t$ is K\"ahler, $P_t$ commutes with $\Delta$; it
therefore preserves $S_t$, and its restriction to $S_t$ is an orthogonal
projection. It suffices to show that $\|P_t\psi\|\le\tfrac12$ for all
unit $\psi\in S_t$ and all small $t$. Suppose not; then there are
$t_l\to0$ and unit $\psi_l\in S_{t_l}$ with $\|P_{t_l}\psi_l\|>\tfrac12$.
Fix $R,A$ and pass to a subsequence as in the preceding proof, so that
$g=g_{t_l}$ is a metric on $\overline X_{\epsilon_l,u_0}$ within relative
distance $\delta\le2\delta(R,A)$ of the model metric $\bar g$. Constants
$C$ below are independent of $R$, $A$ and $l$.

\emph{Comparison with the small eigenforms of the model.} Let
$\mathcal Z$ be the space of closed square-integrable two-forms and
$\bar\psi_l$ the component of $\psi_l$ orthogonal to $\mathcal Z$ for
$\bar g$. Then $d\bar\psi_l=d\psi_l$ and $\bar\psi_l$ is coexact for
$\bar g$. Since $\psi_l$ is orthogonal to $\mathcal Z$ for $g$,
\[
 \|\bar\psi_l\|_{\bar g}=\operatorname{dist}_{\bar g}(\psi_l,\mathcal Z)
 \ge(1-C\delta)\operatorname{dist}_{g}(\psi_l,\mathcal Z)=1-C\delta,
 \qquad\|\psi_l\|_{\bar g}\le1+C\delta,
\]
so $\|\psi_l-\bar\psi_l\|^2_{\bar g}=\|\psi_l\|^2_{\bar g}
-\|\bar\psi_l\|^2_{\bar g}\le C\delta$. Moreover
$\|d\bar\psi_l\|^2_{\bar g}\le(1+C\delta)\lambda^+_{2,m}(t_l)\to0$. By
Proposition~\ref{prop:count} the coexact two-forms of the model have $m$
small eigenvalues and a gap above them, so the distance in $L^2(\bar g)$
from $\bar\psi_l$ to the span $\bar S_l$ of the corresponding eigenforms
tends to zero.

\emph{Limits.} By Proposition~\ref{prop:compactness} and the first step
of the proof of Theorem~\ref{thm:model-law}, after passing to a
subsequence, unit elements of $\bar S_l$ converge, smoothly on compact
subsets of $M_0$ and without loss of $L^2$ norm, to forms $\eta^R_D$
with $D$ orthogonal to $K$ for $G_R$. For such $D$ the form
$\eta^R_D-\eta_D$ lies in the closure of the exact forms in the domain
of $d_{\min}$, to which $\eta_D$ is orthogonal for $g_0$, and $\eta^R_D$
minimizes the norm of $g_0^R$ in its class; with
\eqref{eq:interpolation-errors} this gives
\[
 \|\eta^R_D-\eta_D\|^2_{g_0}=\|\eta^R_D\|^2_{g_0}-\|\eta_D\|^2_{g_0}
 \le CR\,\|\eta_D\|^2_{g_0}.
\]
The forms $\eta_D$ are of type $(1,1)$ on $M_0$ by
Proposition~\ref{prop:harmonic}.

\emph{Conclusion.} For $\rho>0$ let $W_\rho=M_0\cap\{r\ge\rho\}$. The
identifications $\Phi_{t}$ of Lemma~\ref{lem:condition-N} restrict to
embeddings of $W_\rho$ into $X_t$ that depend smoothly on $t$ and are the
inclusion at $t=0$, so the complex structures $\Phi_t^*J_t$ converge
smoothly to that of $M_0$ on $W_\rho$, and so do the projections
$P_t$ to the projection $P_0$ on $M_0$. Combining the two steps, on
$W_\rho$ the forms $\psi_l$ are within $C(\delta^{1/2}+R^{1/2})+o(1)$ in
$L^2$ of a form $\eta_D$ of norm at most $2$, for which $P_0\eta_D=0$;
and the $L^2$ norm of $\psi_l$ on the complement of $W_\rho$ is at most
$C(\delta^{1/2}+R^{1/2})+o(1)+\sup\|\eta_D\|_{L^2(\{r<\rho\})}$, the
supremum over $D$ with $G(D,D)\le4$, which tends to zero with $\rho$.
Since $|P_t\psi|\le|\psi|$ pointwise, choosing first $R$ small and $A$
large, then $\rho$ small, gives $\|P_{t_l}\psi_l\|<\tfrac12$ for large
$l$, a contradiction.
\end{proof}
